\clearpage\pagestyle{otMain}\setlength{\otBlocEspacement}{2.5mm}\fontsize{11}{15}\selectfont\corrigesChapitre{20}{Penser et programmer par blocs}
\begin{blocLivre}{cor-20-1}\begin{corrige}{1}{Entrée et sortie}\demarche{Distinguer donnée saisie et résultat.}
L’entrée est $\res7$. La sortie calculée est $7+5=\res{12}$. Le programme ne dit pas le texte « réponse + 5 » : il évalue l’expression du bloc.
\label{fig-cor-20-1}\figureGeo[corrige]{prog-1}
\end{corrige}\end{blocLivre}
\begin{blocLivre}{cor-20-2}\begin{corrige}{2}{Prévoir une formule}\demarche{Lire les blocs imbriqués.}
Le bloc calcule $3\times4+2=\res{14}$. Le produit se fait avant l’addition.
\label{fig-cor-20-2}\figureGeo[corrige]{prog-2}
\end{corrige}\end{blocLivre}
\begin{blocLivre}{cor-20-3}\begin{corrige}{3}{Des parenthèses}\demarche{Respecter l’emboîtement.}
On calcule d’abord $4+2=6$, puis $3\times6=\res{18}$.
\label{fig-cor-20-3}\figureGeo[corrige]{prog-3}
\end{corrige}\end{blocLivre}
\begin{blocLivre}{cor-20-4}\begin{corrige}{4}{Une aire}\demarche{Interpréter la formule.}
$5\times5=\res{25}$. Si la réponse est un côté en cm, la sortie représente une aire de $\res{25\ \mathrm{cm}^2}$.
\label{fig-cor-20-4}\figureGeo[corrige]{prog-4}
\end{corrige}\end{blocLivre}
\begin{blocLivre}{cor-20-5}\begin{corrige}{5}{Deux déplacements}\demarche{Lire une séquence dans l’ordre.}
Il parcourt $30+50=\res{80}$ pas vers la droite. Les deux instructions sont exécutées successivement.
\label{fig-cor-20-5}\figureGeo[corrige]{prog-5}
\end{corrige}\end{blocLivre}
\begin{blocLivre}{cor-20-6}\begin{corrige}{6}{Modifier un paramètre}\demarche{Modifier le coefficient utile.}
Remplacer 2 par 3. La nouvelle sortie pour 6 vaut $3\times6+4=\res{22}$. L’addition de 4 reste après le produit.
\label{fig-cor-20-6}\figureGeo[corrige]{prog-6}
\end{corrige}\end{blocLivre}
\begin{blocLivre}{cor-20-7}\begin{corrige}{7}{Tracer un carré}\demarche{Relier répétitions et côtés.}
Les quatre côtés sont égaux et les quatre angles sont droits : un $\res{\text{carré}}$. Le périmètre vaut $4\times50=\res{200}$ pas. Le lutin revient au départ et retrouve sa direction.
\label{fig-cor-20-7}\figureGeo[corrige]{prog-7}
\end{corrige}\end{blocLivre}
\begin{blocLivre}{cor-20-8}\begin{corrige}{8}{Un rectangle en boucle}\demarche{Compter les instructions à l’intérieur de la boucle.}
Il trace un $\res{\text{rectangle}}$ de côtés 80 et 40 pas. Le périmètre vaut $2(80+40)=\res{240}$ pas. Chaque tour de boucle trace deux côtés.
\label{fig-cor-20-8}\figureGeo[corrige]{prog-8}
\end{corrige}\end{blocLivre}
\begin{blocLivre}{cor-20-9}\begin{corrige}{9}{Un triangle équilatéral}\demarche{Distinguer le virage de l’angle intérieur.}
Il trace trois segments égaux. Les changements de direction de $120^\circ$ donnent des angles intérieurs de $180-120=60^\circ$. La figure se ferme : c’est un $\res{\text{triangle équilatéral}}$.
\label{fig-cor-20-9}\figureGeo[corrige]{prog-9}
\end{corrige}\end{blocLivre}
\begin{blocLivre}{cor-20-10}\begin{corrige}{10}{L’ordre des blocs}\demarche{Suivre position et direction séparément.}
Le premier finit 40 pas à droite du départ ; le second 40 pas au-dessus. Les orientations finales sont identiques, mais les positions diffèrent. L’ordre des instructions a donc une conséquence.
\label{fig-cor-20-10}\figureGeo[corrige]{prog-10}
\end{corrige}\end{blocLivre}
\begin{blocLivre}{cor-20-11}\begin{corrige}{11}{Un pentagone régulier}\demarche{Relier un tour complet aux virages.}
Les cinq changements de direction totalisent $5\times72=360^\circ$. Les cinq côtés égaux forment ici un pentagone régulier fermé ; son périmètre vaut $\res{200}$ pas.
\label{fig-cor-20-11}\figureGeo[corrige]{prog-11}
\end{corrige}\end{blocLivre}
\begin{blocLivre}{cor-20-12}\begin{corrige}{12}{Une ligne brisée}\demarche{Décrire un motif puis le répéter.}
Chaque répétition ajoute 30 pas vers la droite et 20 vers le haut, puis rétablit la direction de départ. Après trois fois : $\res{90}$ pas à droite et $\res{60}$ vers le haut. Le dessin est un escalier.
\label{fig-cor-20-12}\figureGeo[corrige]{prog-12}
\end{corrige}\end{blocLivre}
\begin{blocLivre}{cor-20-13}\begin{corrige}{13}{Réparer un carré incomplet}\demarche{Compter les côtés réellement tracés.}
Il manque le quatrième côté. Remplacer trois répétitions par $\res4$. Le quatrième déplacement ferme la figure ; le dernier virage rétablit l’orientation initiale.
\label{fig-cor-20-13}\figureGeo[corrige]{prog-13}
\end{corrige}\end{blocLivre}
\begin{blocLivre}{cor-20-14}\begin{corrige}{14}{Une formule de volume}\demarche{Reconnaître un produit répété.}
Les sorties sont $2^3=\res8$ et $3^3=\res{27}$. Elles représentent des volumes en $\res{cm^3}$. Une seule donnée saisie est lue trois fois dans le calcul.
\label{fig-cor-20-14}\figureGeo[corrige]{prog-14}
\end{corrige}\end{blocLivre}
\begin{blocLivre}{cor-20-15}\begin{corrige}{15}{Concevoir puis contrôler}\demarche{Séparer le calcul et le sens géométrique.}
Demander le côté puis dire $4\times\text{réponse}$. Pour 2,5, la sortie est $\res{10\ \mathrm{cm}}$ ; pour 0, elle est 0. Ce test contrôle la formule à la frontière, mais un côté nul ne donne pas un carré non aplati.
\label{fig-cor-20-15}\figureGeo[corrige]{prog-15}
\end{corrige}\end{blocLivre}
